% Transcription — fonds Alexandre Grothendieck, shelfmark 15, batch 7 (pages 121-140).
% Established from the facsimile archives/batches/15/batch-07.pdf (a mirror of
% https://grothendieck.umontpellier.fr/15.pdf), per the skill
% transcribe-grothendieck. The batch file carries no cover sheet: PDF page N
% is page 120+N of the folder (the Montpellier footer prints « 121 » ... « 140 »),
% and the archivists' pencil numbers bottom-left agree leaf by leaf wherever
% legible (121, 122, 127-131, 134, 136, 138, 140 read directly; the others
% are faint but in sequence).
%
% Pass: Opus 5.5 (claude-opus-5-5), 2026-09-23 — first pass, unchecked against the pages by a human.
%
% Eleven of the twenty pages are transcribed. Absent: pages 121, 123, 125,
% three leaves of an unrelated English typescript (a translated paper on
% regular maps in the Zariski topology, its own pages 77, 76, 74, with a
% marginal « p. 210 » referring to the original; nothing in his hand);
% pages 126 and 133, blank; page 135, an administrative typed sheet
% (« DIVERS — CANDIDATS », dated 1.4.68 — the only date in the batch, on a
% verso, so no evidence for the notes themselves); page 137, a typed errata
% list of some other text, scanned upside down; page 139, an orange cover
% inscribed in his hand « Motifs essentiellement abéliens et théorie du
% corps de classes », whose words become the \section heading before page
% 140, per references/hand.md §5.
%
% What the batch is. Four runs in his hand, none paginated by him.
%   (1) Pages 122 and 124, a very fast hand, continuing a sentence begun on
%   page 120 (batch 6): the cocharacters i : G_{m,C} → G_C and the
%   « structure de positivité » (polarisations positives) of the motivic
%   Galois group G; the differences i - i' span a finite-index subgroup of
%   Hom(G_{m,C}, G^0_C), with no indication of the index (« Question semble
%   sérieuse »); R-forms defined by the various i. Then, over a finite field,
%   G_C = D_C(M) with M ⊂ L^* (Weil numbers), complex conjugation inducing a
%   canonical involution u ↦ u' which on Hom(G_{m,C}, G_C) is complex
%   conjugation; page 124 doubts that it lifts to the gerbe, invoking a
%   fundamental class in H^?(Q, U). Much of the connective prose is \ill{}.
%   (2) Pages 127-132, written fairly: Definition (totally real fields and
%   elements, totally positive), Scholie; the maximal totally real field 𝕂
%   ⊂ Q̄ corresponds to the closed normal subgroup π' generated by complex
%   conjugation τ; H = π'^{ab} is killed by 2; the unique quadratic extension
%   𝕃 of 𝕂 Galois over Q, three equivalent characterisations (i)-(iii), with
%   proof; Corollary: L ⊂ 𝕃 iff L is totally real or a totally imaginary
%   quadratic extension of a totally real field (CM); NB: 𝕃 is the largest
%   Galois extension in whose group τ is central, so it contains Q^{ab}.
%   Proposition (130): for Q ∈ Q^* and λ algebraic, (a) |λ_α|^2 = Q for all
%   conjugates ⟺ (b), (b'), (b'') [λ root of λ^2 - aλ + Q with a totally
%   real, 4Q - a^2 totally ≥ 0] ⟺ (c) λ ∈ 𝕃 and |λ|^2 = Q; proof of all
%   implications. A cancelled boxed corollary. Remark (131-132): for K
%   totally real and Q > 1 rational there is a CM quadratic extension of K
%   generated by λ with λλ̄ = Q^i (λ integral if Q ∈ Z), and then λ^{1/i}.
%   (3) Pages 134, 136, 138, headed « Groupes de Galois motiviques
%   commutatifs »: M the category of (semi-simple) isomotives over k, G its
%   (commutative) band, 𝒢 the gerbe of fibre functors; the Tate subcategory
%   M_0 gives 𝒢 → 𝒢_0, ε : G → G_m; a fibre functor φ on M_0 gives a gerbe
%   𝒢_φ under Ker ε; the subcategory M' (Artin-type objects, « Q-vectoriels
%   tordus ») gives G → G' with kernel G^0 and, via k̄, a gerbe 𝒢_ψ under
%   G^0; M^{(0)} generated by M_0 and M' has G^{(0)} = G_m × G', whence a
%   gerbe 𝒢_χ under U = Ker(G → G^{(0)}). Page 138 (a torn sheet written
%   sideways, turned upright here) draws the squares of gerbes and bands and
%   their behaviour under an algebraic extension k'/k: U' ≃ U, G'^0 ≃ G^0.
%   (4) Page 140, after the cover: §1 (circled), T_K = R_{K/Q} G_m, V_K =
%   Ker N_{K/Q}; NB on units; Proposition: for a quotient V'_K = V_K/R,
%   equivalence of (i) finite image of the units, (ii) a modulus m with T →
%   T_m killing R, (iii) (V'_K)_R a compact torus, (iv) R contains the kernel
%   of V_K → V_L → V_L/Im V_{L_0}, L the largest « Weilien » (CM) subfield;
%   the proof of (i)⟺(ii) (Chevalley) and (i)⟹(iii) begins, and the page
%   ends mid-sentence (« E_1 (donc F) est un »), continued presumably in
%   batch 8.
%
% Notation fixed for the batch. \mathcal{G}, \mathcal{G}_0, \mathcal{G}',
% \mathcal{G}_\varphi, \mathcal{G}_\psi, \mathcal{G}_\chi, \mathcal{G}^{(0)} for
% his curly G (gerbes); G, G^0, G', G^{(0)} for the bands (he writes G° for
% the neutral component; G^0 here); \mathcal{M}, \mathcal{M}_0, \mathcal{M}',
% \mathcal{M}^{(0)} for his curly M (categories), M for objects; U, V plain
% for his cursive U, V (U is sometimes a script 𝒰 on the page); \mathbb{K},
% \mathbb{L} for his double-struck K, L (maximal totally real, maximal CM),
% K, L plain where he writes them plain. His « tot. », « tot^t », « ss-corps »,
% « ss-groupe », « hom. », « isom. », « comp. », « ssi » are kept. Plain square
% brackets in the text are his: his own brackets, or words added in the
% interline; editorial additions are \add{}.
%
% Legibility. Pages 127-132 and 140 read almost end to end. Pages 122 and
% 124 are a fast palimpsest whose formulas carry and whose prose largely does
% not; page 134/136 are in between. The rotated marginal notes of 134 and 136
% are partly \ill{}.
%
% The dating is the inventory's own for the folder, copied verbatim. Its
% group, [10-18] « Théorie des motifs », is dated [à partir de 1958-à partir
% de 1982].
\documentclass[11pt,a4paper]{article}
% Le préambule commun à toutes les transcriptions du fonds Grothendieck.
%
% Il tient en une idée : **l'incertitude se compose, elle ne se résout pas.**
% Une transcription qui lisse un mot douteux a détruit la seule chose pour
% laquelle on la faisait — un lecteur doit pouvoir savoir, à chaque endroit,
% ce qui était sur le feuillet et ce qui a été ajouté. Les six macros qui
% suivent sont donc l'appareil critique, et non de la décoration.
%
% Les mêmes commandes sont reconnues par `scripts/render.mjs`, qui produit la
% vue de lecture du site. Une macro ajoutée ici doit l'être aussi là-bas,
% faute de quoi le rendu échouera bruyamment — ce qui est voulu : un rendu
% silencieusement faux est pire qu'un rendu absent.

\ProvidesPackage{grothendieck}[2026/08/08 Transcriptions du fonds Grothendieck]

\usepackage{amsmath,amssymb,amsthm}
\usepackage{mathrsfs}
\usepackage{stmaryrd}
% Les diagrammes commutatifs sont partout dans ce fonds : c'est la langue
% même de la géométrie algébrique de Grothendieck, et une transcription qui
% les rendrait en prose ne transcrirait rien.
\usepackage{tikz-cd}
% Français par défaut — c'est la langue des feuillets — mais l'anglais est
% chargé aussi : la traduction et le résumé sont des documents anglais, et la
% typographie française (espaces avant les deux-points) y serait une faute.
% Ces documents basculent par \selectlanguage{english} en tête de corps.
\usepackage[english,french]{babel}
\usepackage{fontspec}
\usepackage{geometry}
\geometry{a4paper,margin=25mm,marginparwidth=28mm}
\usepackage{marginnote}
\usepackage{xcolor}
% ulem, not soul. `soul`'s \st and \ul are built on a character-by-character
% scan that XeLaTeX's font machinery breaks: the document fails with an
% undefined control sequence rather than degrading. ulem does the same two jobs
% and is Unicode-engine safe. `normalem` stops it hijacking \emph, which the
% transcriptions use for Grothendieck's own emphasis.
\usepackage[normalem]{ulem}
\usepackage{hyperref}
\hypersetup{colorlinks=true,linkcolor=black,urlcolor=black,pdfborder={0 0 0}}

% --- Métadonnées du lot -----------------------------------------------------
% Reprises telles quelles par le script de rendu, qui les affiche en tête de
% la vue de lecture. Elles fixent ce que le fichier prétend couvrir : sans
% elles, une transcription flotte sans point d'attache dans un fonds de seize
% mille feuillets.

\newcommand{\folder}[1]{\gdef\grothFolder{#1}}
\newcommand{\batch}[1]{\gdef\grothBatch{#1}}
\newcommand{\leaves}[2]{\gdef\grothFirst{#1}\gdef\grothLast{#2}}
\newcommand{\foldertitle}[1]{\gdef\grothFoldertitle{#1}}
\gdef\grothFolder{?}\gdef\grothBatch{?}\gdef\grothFirst{?}\gdef\grothLast{?}\gdef\grothFoldertitle{}

% --- Le résumé --------------------------------------------------------------
%
% Il ouvre la lecture modernisée, sous le titre « Résumé », et il est composé
% en petit. Ce n'est pas un ornement : c'est le seul passage du document qui
% ne soit pas des mathématiques, et un lecteur doit pouvoir le reconnaître
% avant de l'avoir lu — puis le sauter s'il connaît déjà le sujet, ou s'y
% arrêter s'il ne le connaît pas. Le filet qui le suit marque la reprise du
% corps.

\newenvironment{resume}
  {\small\setlength{\parskip}{0.45em}}
  {\par\medskip\hrule\medskip}

% --- Mots-clés ---------------------------------------------------------------
%
% En anglais, en clôture du résumé de la lecture modernisée : le vocabulaire
% moderne sous lequel chercher ce que les feuillets construisent. Le manifeste
% du site (`npm run manifest`) les extrait pour étiqueter la cote entière —
% cette ligne est la source unique des tags, il n'y a pas de fichier à côté.
\newcommand{\keywords}[1]{\par\smallskip\noindent
  {\footnotesize\textsc{Keywords} --- \textit{#1}}\par}

% --- L'appareil critique ----------------------------------------------------

% Le numéro de feuillet, en marge. C'est l'ancre qui permet de régler un
% désaccord sur une formule en regardant *un* feuillet plutôt qu'une chemise
% de six cents. Le site s'en sert aussi pour tourner les pages du fac-similé
% quand on fait défiler la transcription.
\newcommand{\leaf}[1]{%
  \marginnote{\footnotesize\color{gray!70}\textsf{#1}}%
  \label{leaf:#1}\ignorespaces}

% Illisible : on ne devine pas. Un mot inventé qui se lit comme les autres est
% le pire résultat possible.
\newcommand{\ill}{\textcolor{red!60!black}{[\,\ldots\,]}}

% Lecture douteuse : le mot est donné, et signalé comme incertain.
\newcommand{\uncertain}[1]{\uline{#1}}

% Ajout de l'éditeur — une lettre manquante, un mot sous-entendu.
\newcommand{\add}[1]{\textcolor{blue!50!black}{[#1]}}

% Biffé par Grothendieck lui-même. Ce qu'il a rayé fait partie du document :
% une rature dit souvent où la pensée a changé de direction.
\newcommand{\struck}[1]{\sout{#1}}

% Note du transcripteur, sur l'état matériel du feuillet.
\newcommand{\note}[1]{\par\smallskip{\small\color{gray!60!black}\textit{[#1]}}\par\smallskip}

% Note marginale de Grothendieck, distincte du corps du texte.
\newcommand{\marginal}[1]{\par\smallskip\noindent\hspace{1em}%
  {\small\color{gray!60!black}#1}\par\smallskip}

% --- Titre ------------------------------------------------------------------

\AtBeginDocument{%
  \begin{center}
    {\large\bfseries Fonds Alexandre Grothendieck --- cote n\textsuperscript{o}~\grothFolder}\\[2pt]
    {\itshape\grothFoldertitle}\\[4pt]
    {\small feuillets \grothFirst--\grothLast\ (lot \grothBatch)}\\[2pt]
    {\footnotesize\color{gray!60!black}Transcription établie d'après le fac-similé numérisé
     par l'Université de Montpellier.\\
     Les lectures incertaines sont soulignées, les illisibles notées [\,\ldots\,],
     les ajouts entre crochets.}
  \end{center}
  \vspace{1em}}

\endinput


\folder{15}
\batch{7}
\pages{121}{140}
\dating{1965-[vers 1977]}
\watermark{Édition de démonstration}
\foldertitle{Théorie arithmétique. Théorie de Galois des motifs (1965) : notes manuscites (s.d.), tapuscrit (s.d.).}

\begin{document}

\page{122}
\note{la page s'ouvre au milieu d'une phrase commencée au feuillet 120
(lot 6)}
de $i$. \struck{Question de \uncertain{structure} de positivité et \ill{}}
\struck{\uncertain{analogie} \ill{} $i$ \uncertain{fixé}}
\note{deux lignes biffées, la première d'un trait, la seconde d'un zigzag
encadré}
Donc, \uncertain{grâce} : en \uncertain{toutes} circonstances
\uncertain{spéciales}, il y a des \ill{} dans la \struck{\uncertain{positivité}}
structure (polarisations positives) du $G$. Les $i$ \uncertain{fixés} avec $i$
dans la classe ci-dessous.
\[
  \mathbb{G}_{m,\mathbb{C}} \xrightarrow{\; i \;} G_{\mathbb{C}} ,
  \qquad
  \mathbb{G}_{m,\mathbb{C}} \xrightarrow{\; i' \;} G_{\mathbb{C}}
\]
\note{ces deux flèches sont écrites dans la marge de gauche, en regard des
lignes qui les commentent}
Quid de l'arbitraire ? On constate que l'ens.\ [des] $i - i'$
\struck{conjugués : \ill{}$(\lambda)$, \ill{}$(\lambda)$} définit un
\uncertain{s-groupe} d'indice fini de $\mathrm{Hom}(\mathbb{G}_{m,\mathbb{C}},
G^{0}_{\mathbb{C}})$, mais cela ne donne aucune indication sur l'indice,
\uncertain{ou} désigner la composante des $i_{0} - i'_{0}$ \ill{}
\uncertain{contenant} $i$ fixé. Question semble sérieuse.

\note{le paragraphe suivant est marqué d'un double trait vertical en marge}
En tous cas, \ill{} un \struck{deux} \uncertain{groupe} de $\mathbb{R}$-formes
\uncertain{fibres} \uncertain{polarisées}, définies au moyen des différents
$i$.

Autre structure spécifique sur le $G$, pour \uncertain{$k$} un corps fini :
d'après la \uncertain{forme} de \ill{}, $G_{\mathbb{C}} =
D_{\mathbb{C}}(M)$, $M \subset L^{*}$, comme $L$ est \ill{}, tel que la
\uncertain{conjugaison} complexe, qui commute \add{à}
$\mathrm{Gal}(L/\mathbb{Q})$, \uncertain{induit} sur $M$ donc \add{sur} $G$ une
involution canonique, notée $u \mapsto u'$. Dans \struck{\ill{}} \ill{}
n'est \uncertain{pas} \uncertain{spécialisée}, d'ailleurs \uncertain{fait} sur
\ill{} \ill{} l'identité, ce qui suffit pratiquement, la \ill{}. Sur
$\mathrm{Hom}(\mathbb{G}_{m,\mathbb{C}}, G_{\mathbb{C}})$, $u \mapsto u'$
induit la conjugaison complexe.

\page{124}
Il ne semble pas que cette involution corresponde à une involution sur le
\struck{groupe} $\mathcal{G}$ associé : $M$ (en \struck{\ill{}}
\struck{\uncertain{tenant}} \ill{} \uncertain{de} la propriété de
\uncertain{connexion} \ill{} \uncertain{par} \ill{} l'élément fondamental
\uncertain{composante} des $H^{2}(\mathbb{Q}, \struck{\ill{}}\,U)$ !, qui peut
\struck{\uncertain{d'autres}} \ill{} d'ordre \uncertain{divisé par} $2$ …).
\note{l'exposant de $H$ est mal formé et lu $2$ sans certitude ; devant $U$,
une lettre noircie. La page s'arrête ici, au premier tiers}

\page{127}
\emph{Définition}. Soit $K/\mathbb{Q}$ une extension [algébrique]
\struck{\uncertain{de degré fini}}. On dit que $K$ est tot.\ réel si pour tout
$\mathbb{Q}$-homom.\ $\sigma : K \to \mathbb{C}$, on a $\sigma(K) \subset
\mathbb{R}$. On dit qu'un élément $x$ de $K$ est tot.\ réel [(resp.\ tot.\
$> 0$, tot.\ $\geqslant 0$)] si pour tout $\sigma$ comme dessus, on a
$\sigma(x) \in \mathbb{R}$ (resp.\ $\sigma(x) > 0$, resp.\ $\sigma(x)
\geqslant 0$). Tot$^{t}$ (ou purement) imaginaire : ….

\emph{Scholie}. Pour que $x$ soit tot.\ réel, il faut et il suffit que
$\mathbb{Q}(x) \subset K$ le soit. Pour qu'il soit tot.\ $\geqslant 0$, il
faut et il suffit qu'il soit nul ou tot $> 0$. Les éléments tot $\geqslant 0$
\struck{\ill{}} [sont] les éléments $\geqslant 0$ pour une relation d'ordre
sur $K$. [On \uncertain{notera} \ill{} $x$ tot.\ $\geqslant y$ ou $x$ tot.\
$\leqslant y$]. Si $K$ \struck{est} [est] engendré par des ss-corps
$K_{\alpha}$, alors $K$ est tot.\ réel ssi les $K_{\alpha}$ le sont. Si $K$
est engendré par des $x_{\alpha}$, $K$ tot.\ réel ssi \struck{\ill{}} les
$x_{\alpha}$ le sont.

Il existe, pour $K$ donné, un plus grand ss-corps totalement réel de $K$,
appelé la \emph{clôture tot.\ réelle} de $\mathbb{Q}$ dans $K$. Soit
\struck{$K$} $\overline{\mathbb{Q}} \subset \mathbb{C}$, et $\pi =
\mathrm{Gal}(\overline{\mathbb{Q}}/\mathbb{Q})$, \struck{\ill{}} $\tau \in
\pi$ la conjugaison complexe
\[
  \tau\lambda = \overline{\lambda} \qquad \text{pour } \lambda \in
  \overline{\mathbb{Q}}
\]
de sorte que
\[
  \tau^{2} = 1 .
\]
Alors le ss-groupe $\pi'$ de $\pi$ qui \struck{invariant} [correspond] à la
clôture tot.\ réelle \struck{$K$} $\mathbb{K}$ \struck{\ill{}} de
$\mathbb{Q}$ dans $\overline{\mathbb{Q}}$ est le \struck{\ill{}} ss-groupe
\emph{invariant fermé} engendré par $\tau$, i.e.\ le ss-groupe invariant
engendré par les ${}^{\sigma}\tau = \sigma\tau\sigma^{-1}$ ($\sigma \in \pi$).

\marginal{dans la marge de gauche, une tour verticale : $A$, $L$, $K$,
$\mathbb{Q}$, reliés par des traits}
Considérons la plus grande extension abélienne $A$ de $\mathbb{K}$, dont le
groupe de Galois est donc $H = \pi'/(\pi', \pi')$. \struck{Son} Sur
\note{la phrase se poursuit au feuillet suivant ; le complément de « Sur »
(sans doute $H$) n'est pas écrit}

\page{128}
le groupe $G = \pi/\pi' = \mathrm{Gal}(\mathbb{K}/\mathbb{Q})$ opère
canoniquement, et il est engendré topologiquement par les
${}^{\sigma}\dot{\tau}$, où $\dot{\tau} \in H$ est l'image de $\tau$, et
$\sigma \in G$. Donc il est annulé par $2$.

Soit $L$ une extension abélienne de $\mathbb{K}$, \struck{correspondant à la
\ill{}} \struck{les conditions} [son] groupe de Galois est donc un quotient
$\Gamma$ de $H$ ; \struck{conditions} supposons $L \neq \mathbb{K}$, i.e.\
$\Gamma \neq \{e\}$. Conditions équivalentes :

(i) les $\sigma\tau\sigma^{-1}$ ($\sigma \in \pi$) coïncident toutes sur $L$
(donc y coïncident avec $\tau : \lambda \mapsto \overline{\lambda}$)

(ii) $L$ est une extension \emph{quadratique} de $\mathbb{K}$ \emph{galoisienne
sur} $\mathbb{Q}$.

(iii) $L$ est une extension \emph{quadratique} \emph{purement imaginaire} de
$\mathbb{K}$. \struck{\ill{}}
\note{(iii) est écrit en interligne au-dessus d'une ligne raturée en zigzag}

De plus, une telle $L$ \struck{existe et} est unique, [et d'ailleurs existe :
il suffit de prendre une extension quadratique imaginaire $E$ de $\mathbb{Q}$
(p.\ ex.\ $\mathbb{Q}(\sqrt{-1})$), et $L = \mathbb{Q}(\mathbb{K}, E)$ (p.\
ex.\ $L = \mathbb{K}(\sqrt{-1})$)].
\note{le paragraphe des conditions (i) à (iii) et de l'unicité est encadré à
gauche d'un grand crochet}

\emph{Dém.} (i) $\Rightarrow$ (ii). Comme $\Gamma$ est engendré par les images
\struck{\ill{}} des ${}^{\sigma}\tau$, \struck{\ill{}} et que l'hyp.\ de (i)
signifie que ces images sont égales, on voit que ces images, que $\Gamma$ est
engendré par l'élément $\dot{\tau}$ d'ordre $\mid 2$, donc est $\simeq
\mathbb{Z}/2\mathbb{Z}$ (car $\Gamma \neq (e)$ par hypothèse). D'ailleurs, si
un tel $\Gamma$ existe, \struck{\ill{}} [il est] nécessairement égal à $H_{G}$
(car \struck{\ill{}} \ill{} est un quotient de $H_{G}$, et l'argument
précédent montre que \struck{\ill{}} $H_{G}$ est engendré par un élément
d'ordre $\mid 2$. \struck{\ill{}} \struck{il existe tel et} Il en résulte
aussitôt que $L$ est galoisien sur $\mathbb{Q}$, donc (i) $\Rightarrow$ (ii).
D'autre part, (ii) $\Rightarrow$ (i), car (ii) implique que $\Gamma \simeq
\mathbb{Z}/2\mathbb{Z}$ [et que $G$ y opère], et comme $\Gamma$ est engendré
par les images des $\sigma\tau\sigma^{-1} = {}^{\sigma}\tau$, i.e.\ par les
${}^{\sigma}\dot{\tau}$, on voit que $\dot{\tau}$ (donc les
${}^{\sigma}\dot{\tau}$) sont tous égaux : l'unique élément \ill{}
\struck{\ill{}}. (i), (ii) $\Longleftrightarrow$ (iii) car on remarque que $L$
est quadratique sur $\mathbb{K}$, la condition (i) sur les
$\sigma\tau\sigma^{-1}$ est évidemment équivalente à (iii).

\page{129}
\emph{Corollaire}. Soit $L$ un \struck{\ill{}} sous-corps de
$\overline{\mathbb{Q}}$. Pour qu'on ait $L \subset \mathbb{L}$, il faut et il
suffit que $L$ soit tot.\ réel, ou extension quadratique purement imaginaire
d'un corps tot.\ réel.

\begin{tikzcd}[column sep=small, row sep=small]
\mathbb{K} \arrow[r, no head] \arrow[d, no head] & \mathbb{K}(L) \arrow[r, no head] \arrow[d, no head] & \mathbb{L} \\
K = \mathbb{K} \cap L \arrow[r, no head] \arrow[d, no head] & L & \\
\mathbb{Q} & &
\end{tikzcd}

\note{ce schéma est dans la marge de gauche}

Suffisance triviale par (iii). Nécessité : Par théorie de Galois, $L$ est une
extension de degré $1$ ou $2$ du corps tot.\ réel $K = L \cap \mathbb{K}$, et il
reste à prouver que lorsque $L$ en est une extension quadratique, cette
dernière est tot$^{t}$ imaginaire. \struck{\ill{}} On a \struck{\ill{}} on est
ramené à prouver : \struck{\emph{Lemme}. Soit $L$ une extension quadratique
d'un corps tot.\ réel $K$. Si $L$} (quitte à remplacer $L$ par un conjugué,
[qui est engendré par $K$]) \ill{} $L$ \ill{} $\mathbb{K}(L)$ est imaginaire.
Or s'il était réel, comme $\mathbb{K}$ l'est, \ill{} mais, on a
$\mathbb{K}(L) = \mathbb{L}$, absurde.

\emph{NB}. $\mathbb{L}$ est la plus grande extension galoisienne de
$\mathbb{Q}$ telle que dans son groupe de Galois $G'$, l'image de $\tau$
(conjugaison complexe) soit \emph{centrale}. Donc $\mathbb{L}$ contient
l'extension \uncertain{abélienne} maximale de $\mathbb{Q}$, i.e.\ le corps des
racines de l'unité.

\page{130}
\emph{Proposition}. Soit $Q \in \mathbb{Q}^{*}$, et soit $\lambda \in \mathbb{C}$
un nombre algébrique. Conditions équivalentes :

(a) Pour tout conjugué $\lambda_{\alpha}$ de $\lambda$, on a
$\lambda_{\alpha}\overline{\lambda}_{\alpha} = Q$ i.e.\ $|\lambda_{\alpha}|^{2}
= Q$.

(b) \struck{\ill{}} Il existe un corps [de nbres] \emph{totalement réel} $K$,
tel que \struck{\ill{}} $K(\lambda)$ soit égal à $K$ ou une extension
\emph{quadratique} purement imaginaire de $K$, et on a
$\lambda\overline{\lambda} = Q$ i.e.\ $|\lambda|^{2} = Q$.

(b$'$) Comme (b), avec $K = \mathbb{Q}(a)$ \struck{\ill{}}, où $a = \lambda +
\overline{\lambda} = \lambda + \frac{Q}{\lambda}$.

(b$''$) Il existe un nombre algébrique \struck{$a$} \emph{tot.\ réel} $a$,
avec \struck{\ill{}} $4Q - a^{2}$ nul ou totalement positif, tel que $\lambda$
satisfasse une équation
\[
  \lambda^{2} - a\lambda + Q = 0 .
\]

(c) Soit $\mathbb{K}$ l'extension alg.\ tot.\ réelle maximale de $\mathbb{Q}$,
et $\mathbb{L}$ l'unique extension quadratique de $\mathbb{K}$ galoisienne sur
$\mathbb{Q}$. On a $\lambda \in \mathbb{L}$, et \struck{\ill{}}
$\lambda\overline{\lambda} = Q$ i.e.\ $|\lambda|^{2} = Q$.

\emph{Dém.}
a) $\Longrightarrow$ b$''$) Posons $a = \lambda + \overline{\lambda} = \lambda +
\frac{Q}{\lambda}$. Je dis que $a$ \emph{est tot.\ réel}. En effet, les
conjugués de $a$ sont les \struck{\ill{}} $a_{\alpha} = \lambda_{\alpha} +
\frac{Q}{\lambda_{\alpha}}$, or par hypothèse
$\lambda_{\alpha}\overline{\lambda}_{\alpha} = Q$ donc $a_{\alpha} =
\lambda_{\alpha} + \overline{\lambda}_{\alpha} \in \mathbb{R}$. D'autre part,
$\lambda$ est solution de l'équation
\[
  (T - \lambda)(T - \overline{\lambda}) = 0
  \quad \text{i.e.} \quad
  T^{2} - aT + Q = 0 ,
\]
et par suite $\lambda_{\alpha}$ est solution de l'équation $T^{2} -
a_{\alpha}T + Q = 0$, \struck{de plus} l'autre solution étant
$\frac{Q}{\lambda_{\alpha}}$ [qui est égale : $\overline{\lambda}_{\alpha}$]
\struck{\ill{}} comme \ill{} $\lambda_{\alpha}\overline{\lambda}_{\alpha} =
Q$, les deux solutions sont $\lambda_{\alpha}$ et $\overline{\lambda}_{\alpha}$,
donc complexes conjuguées. Cela montre que $a_{\alpha}^{2} - 4Q \leqslant 0$,
\struck{d'où} i.e.\ $4Q - a_{\alpha}^{2} \geqslant 0$, donc $4Q - a^{2}$ est
tot.\ positif].
\note{sur la page, le coefficient constant de la seconde équation porte un
indice raturé}

b$''$) $\Longrightarrow$ b$'$) Il est \uncertain{évident} que le corps
$\mathbb{Q}(a)$ est tot.\ réel, et $\mathbb{Q}(\lambda) \supset \mathbb{Q}(a)$,
[$\lambda$ est de degré au plus deux sur $\mathbb{Q}(a)$] [car $a = \lambda +
\frac{Q}{\lambda}$]. \struck{\ill{}} \struck{De plus} \struck{\ill{}} Comme
$a^{2} - 4Q \leqslant 0$, les solutions de l'équation $T^{2} - aT + Q = 0$
sont $\lambda$ et $\overline{\lambda}$ (cf.\ \uncertain{comme plus bas}), et
on a $\lambda\overline{\lambda} = Q$, $\lambda + \overline{\lambda} = a$.

\page{131}
[Évident]
\note{« Évident » est écrit au-dessus de « Il faut prouver », qui n'est pas
biffé}
Il faut prouver que si $\lambda \notin \mathbb{K}$ i.e.\ $a^{2} - 4Q \neq 0$
(cf.\ lemme plus bas), alors \struck{$\lambda$ est purement}
$\mathbb{Q}(\lambda)$ est purement imaginaire, i.e.\ les $\lambda_{\alpha}$
sont tous non réels. \struck{Cela signifie que les} Comme $\lambda_{\alpha}$
est solution de l'équation $T^{2} - a_{\alpha}T + Q = 0$, cela signifie que
$a_{\alpha}^{2} - 4Q < 0$, ce qui résulte du fait que $4Q - a^{2}$ est tot
$\geqslant 0$.

b$'$) $\Longrightarrow$ b) trivial, \struck{\ill{}}

b) $\Longrightarrow$ c) trivial, car $\mathbb{K}(\lambda)$ est égal à
$\mathbb{K}$ ou une ext.\ quadratique purement imaginaire de $\mathbb{K}$.

c) $\Longrightarrow$ a) car si $\lambda_{\alpha} = \sigma_{\alpha}(\lambda)$,
on aura [$\lambda_{\alpha} \in \mathbb{L}$, donc] \struck{\ill{}}
$\overline{\lambda}_{\alpha}$ est l'unique \struck{\ill{}} conjugué de
$\lambda_{\alpha}$ sur $\mathbb{K}$, \struck{\ill{}} donc
\[
  \lambda_{\alpha}\overline{\lambda}_{\alpha}
  = N_{\mathbb{L}/\mathbb{K}}(\lambda_{\alpha})
  = N_{\mathbb{L}/\mathbb{K}}(\lambda) = \lambda\overline{\lambda} = Q .
\]

\struck{Supposons $Q \in \mathbb{Q}^{*+}$, $L$ fini sur $\mathbb{Q}$, \ill{}
réel (i.e.\ \ill{} tot.\ réel). \emph{Corollaire} (Soit $L \subset
\mathbb{L}$. Alors $L$ est engendré par un $\lambda$ tel que
$\lambda\overline{\lambda} = Q$. Si $Q \in \mathbb{Z}$, on peut supposer que
$\lambda$ est entier.}
\note{passage encadré et biffé de grandes boucles}

\emph{Remarque}. Soit $K$ une extension finie tot.\ réelle de $\mathbb{Q}$, et
$Q \in \mathbb{Q}^{*+}$, $Q > 1$. Je dis qu'il existe alors une extension
quadr.\ purement imaginaire $L$ de $K$, \struck{telle que} et un entier $i
\geqslant 0$, \struck{tels que} tels que $L$ soit engendré [sur $\mathbb{Q}$]
par un élément $\lambda$ tel que $\lambda\overline{\lambda} = Q^{i}$ i.e.\
$|\lambda|^{2} = Q^{i}$. Si $Q \in \mathbb{Z}$, on peut même prendre $\lambda$
entier algébrique [on peut supposer $a$ entier].

\emph{Dém.} Soit $a \in K$ tel que \struck{$K$} $K = \mathbb{Q}(a)$. [Soit $i$
tel que] \struck{\ill{}} $4Q^{i} - a^{2}$ soit tot.\ $> 0$ ($i$ assez grand
pour que \ill{} ; \struck{\ill{}} existe, car $Q \in \mathbb{Q}$, $Q > 1$).
Soit $\lambda$ solution de l'équation $\lambda^{2} - a\lambda + Q^{i} = 0$.
Il satisfait à la condition voulue ; si $Q$ est entier, il est entier, $a$
l'étant.

\page{132}
Notons que $\lambda' = \lambda^{1/i}$ sera aussi dans $\mathbb{L}$ (car
$\lambda\overline{\lambda} = Q$ \struck{\ill{}}, et de même pour ses
conjugués …), donc on voit ainsi qu'il existe un $\lambda'$ \struck{\ill{}}
$\in \mathbb{L}$, avec $\lambda'\overline{\lambda}' = |\lambda'|^{2} = Q$, tel
que $\mathbb{Q}(\lambda')$ soit \struck{p} (purement) imaginaire et
\emph{contienne} $K$ ; de plus, si $Q$ est entier, on peut prendre $\lambda'$
entier.
\note{« dans $\mathbb{L}$ » : la double barre de la lettre n'est pas nette sur
la page ; le sens l'impose. De même, $\lambda\overline{\lambda} = Q$ est écrit
ainsi, pour $Q^{i}$}

\section{Groupes de Galois motiviques commutatifs}

\note{titre de sa main, souligné, en tête du feuillet 134 ; la rédaction se
poursuit aux feuillets 136 et 138}

\page{134}
$k$ corps

$\mathcal{M}$ catégorie des (iso)motifs (½simples) sur $k$

$G$ lien de $\mathcal{M}$, \struck{supposé commutatif}.

$\mathcal{G}$ gerbe des \struck{\ill{}} foncteurs fibres de $\mathcal{M}$

\struck{$\mathcal{G}$}

\struck{$\mathcal{G}_{0}$} $\to$ $\mathcal{M}_{0} \subset \mathcal{M}$,
catégorie \struck{des Tate} [engendrée] par $\mathbb{Q}(1)$ de Tate

$\mathcal{G}_{0}$, $G_{0} = \mathbb{G}_{m}$ gerbe et lien correspondants
\[
  \mathcal{G} \to \mathcal{G}_{0} ,
  \qquad
  G \xrightarrow{\;\varepsilon\;} \mathbb{G}_{m}
  \qquad \text{épimorphismes}
\]
$\mathcal{G}_{0}$ est \uncertain{neutralisée} i.e.\ munie d'une
\uncertain{section}, correspondant au foncteur fibre $\varphi$
\struck{$\mathbb{Q}(1) \to \mathbb{Q}$} choisi sur $\mathcal{M}_{0}$ [\emph{NB}
Comme $H^{1}(\mathbb{Q}, \mathbb{G}_{m}) = 0$, deux foncteurs fibres sur
$\mathcal{M}_{0}$ sont [\ill{}] isomorphes].

Une trivialisation de $\mathcal{G}_{0}$ permet de \uncertain{restreindre} le
lien de $\mathcal{G}$ de $G$ à : $\operatorname{Ker}\varepsilon = V$ : on
prend les foncteurs fibres \emph{munis} d'un isom.\ de leur restriction à
$\mathcal{M}_{0}$ avec le foncteur fibre fixé sur $\mathcal{M}_{0}$.
\uncertain{Leur} \ill{} une nouvelle gerbe $\mathcal{G}_{\varphi}$ de lien
$U$.
\note{la lettre après « Ker $\varepsilon$ = » est un $V$ écrit en surcharge,
d'une encre plus noire, sur un $U$ ; le $U$ de la fin de l'alinéa n'est pas
corrigé. Au feuillet 138, $V \to G \xrightarrow{\varepsilon}
\mathbb{G}_{m}$}
\marginal{\uncertain{vérif} Giraud ?}

Soit $\mathcal{M}'$ la ss-catégorie de $\mathcal{M}$ formée des objets $M$
tels que les $\otimes(M + \check{M})$ soient combinaisons \uncertain{d'un} nb
fini d'objets simples. [$\mathcal{G}'$ la gerbe associée, $G'$ le lien.]
(\ill{} \ill{} \ill{})
\[
  \begin{array}{l}
    \mathcal{G} \longrightarrow \mathcal{G}' \longrightarrow 1 \\[2pt]
    1 \longrightarrow G^{0} \longrightarrow G
      \xrightarrow{\ \text{épim.}\ } G' \longrightarrow 1
  \end{array}
\]
Le noyau de $G \to G'$ est la « composante neutre » $G^{0}$. D'autre part,
l'on sait que $\mathcal{M}'$ est équivalente à une ss-catégorie [pleine de
celle] des $\mathbb{Q}$-\uncertain{vectoriels} \uncertain{tordus} sur $k$. Le
\uncertain{choix} d'une clôture algébrique $\bar{k}$ de $k$ fournit un

\page{136}
foncteur fibre [$\psi = \psi_{\bar{k}}$] sur $\mathcal{M}'$, i.e.\ une section
de $\mathcal{G}'$, d'où une \struck{\ill{}} \struck{restriction} gerbe
$\mathcal{G}_{\psi}$ de lien $G^{0}$, qui \struck{\ill{}} \struck{\ill{}}
\uncertain{est reliée} à la gerbe $\mathcal{G}$ de lien $G$ via $G^{0} \to G$.

\emph{NB} En général, $H^{1}(\mathbb{Q}, \Gamma) \neq 1$, donc il y a ici un
choix \uncertain{bis}, \ill{} de trivialisation de la gerbe $\mathcal{M}'$.

\marginal{\ill{} \ill{} dans $\mathcal{M}$}
\note{note entourée, écrite en biais dans la marge de gauche}
Considérons enfin la \struck{gerbe} [ss-catégorie \uncertain{engendrée}
$\mathcal{M}^{(0)}$ de $\mathcal{M}$] \struck{\ill{}} [par] $\mathcal{M}_{0}$
et $\mathcal{M}'$ : c'est la catégorie \uncertain{formée} des sommes directes
\[
  \sum_{i \in \mathbb{Z}} V_{i}(i) , \qquad \text{où les}
\]
$V_{i}$ sont $\in \mathcal{M}'$ et les $V_{i} \in \mathcal{M}'$ sont tous nuls
sauf un nb fini [\emph{NB}
\[
  \Bigl(\sum V_{i}(i)\Bigr) \otimes \Bigl(\sum W_{j}(j)\Bigr)
  = \sum_{h} \Bigl(\sum_{i+j=h} V_{i} \otimes W_{j}\Bigr)(h) ,
  \qquad
  \Bigl(\sum V_{i}(i)\Bigr)^{\vee} = \sum \check{V}_{i}(-i) \,\Bigr]
\]
On trouve un foncteur fibre
\[
  \sum_{i \in \mathbb{Z}} V_{i}(i) \longmapsto \sum V_{i}(\bar{k}) ,
\]
d'où $\mathcal{G}^{(0)}$, $G^{(0)} \simeq G' \times \mathbb{G}_{m}$
\note{sous $G^{(0)}$, un mot court, \ill{}. Dans la marge de gauche :}

\begin{tikzcd}[column sep=small, row sep=small]
\mathcal{M}_{0} \arrow[r, no head] & \mathcal{M}^{(0)} \\
 & \mathcal{M}' \arrow[u]
\end{tikzcd}

\[
  \begin{array}{l}
    \mathcal{G} \longrightarrow \mathcal{G}^{(0)} \\[2pt]
    1 \longrightarrow U \longrightarrow G \longrightarrow G^{(0)} =
      \mathbb{G}_{m} \times G' \longrightarrow 1
  \end{array}
\]
On trouve une \struck{neutralisation} [gerbe] $\mathcal{G}_{\chi}$ sous $U$,
reliée à la gerbe $\mathcal{G}$ sous $G$ \struck{\ill{}} et $U
\hookrightarrow G$.

\page{138}
\note{feuillet déchiré, écrit perpendiculairement à la cote ; on le lit ici
tourné d'un quart de tour}

\begin{tikzcd}
\mathcal{G}_{\varphi} \arrow[r] & \mathcal{G} \\
\mathcal{G}_{\chi} \arrow[u] \arrow[r] & \mathcal{G}_{\psi} \arrow[u]
\end{tikzcd}

\begin{tikzcd}
 & G' & \\
V \arrow[r] & G \arrow[u] \arrow[r, "\varepsilon"] & \mathbb{G}_{m} \\
U \arrow[u] \arrow[r] & G^{0} \arrow[u] &
\end{tikzcd}

Soit $k'$ une extension algébrique de $k$, d'où un foncteur
\[
  \begin{array}{l}
    \mathcal{M} \longrightarrow \mathcal{M}' \quad \text{catégorie des motifs
      sur } k' \text{ engendrée par les } M \otimes_{k} k',\ M \in \mathcal{M}
      . \\[2pt]
    \mathcal{M}^{(0)} \longrightarrow \mathcal{M}'^{(0)}
  \end{array}
\]
d'où
\[
  \mathcal{G}' \longrightarrow \mathcal{G}
\]

\begin{tikzcd}
G' \arrow[r, hook] \arrow[d] & G \arrow[d] \\
\Gamma' \arrow[r, no head] & \Gamma
\end{tikzcd}

\note{sous $\Gamma'$ et $\Gamma$, reliés par un signe d'égalité vertical :
$G'/G'^{0}$ et $G/G^{0}$. Ici les accents désignent les objets relatifs à
$k'$, non plus ceux de $\mathcal{M}'$ des feuillets 134-136}

L'homom.\ $G' \hookrightarrow G$ \struck{définit} est un \uncertain{monom.}\
\ill{}, définit un isom.\ sur les comp.\ neutres.

\begin{tikzcd}
\mathcal{G}'_{\varphi'} \arrow[r, no head] & \mathcal{G}' \\
\mathcal{G}'_{\chi'} \arrow[u] \arrow[r] & \mathcal{G}'_{\psi'} \arrow[u]
\end{tikzcd}

\begin{tikzcd}
\mathcal{G}_{\varphi} \arrow[r, no head] & \mathcal{G} \\
\mathcal{G}_{\chi} \arrow[u] \arrow[r, no head] & \mathcal{G}_{\psi} \arrow[u]
\end{tikzcd}

\note{sur la page, une grosse flèche horizontale va du premier carré au
second}

\begin{tikzcd}
V' \arrow[r] & G' \\
U' \arrow[u] \arrow[r] & G'^{0} \arrow[u]
\end{tikzcd}

\begin{tikzcd}
V \arrow[r, no head] & G \\
U \arrow[u] \arrow[r, no head] & G^{0} \arrow[u]
\end{tikzcd}

\note{de même, une grosse flèche horizontale va du premier carré au second}

En particulier, on trouve
\[
  \mathcal{G}'_{\chi'} \longrightarrow \mathcal{G}_{\chi} ,
  \qquad
  \mathcal{G}'_{\psi'} \longrightarrow \mathcal{G}_{\psi} ,
\]
\[
  U' \xrightarrow{\;\sim\;} U ,
  \qquad
  G'^{0} \simeq G^{0} .
\]

\section{Motifs essentiellement abéliens et théorie du corps de classes}

\note{titre de sa main sur une couverture orange (feuillet 139), qui ne
porte rien d'autre}

\page{140}
\marginal{§1 (entouré)}
$K$ extension finie de $\mathbb{Q}$
\[
  T_{K} = \prod_{K/\mathbb{Q}} \mathbb{G}_{m} ,
  \qquad
  N_{K/\mathbb{Q}} : T_{K} \to T_{\mathbb{Q}} = \mathbb{G}_{m} ,
  \qquad
  V_{K} \, (\subset T_{K}) = \operatorname{Ker} N_{K/\mathbb{Q}} ;
\]
\note{devant $V_{K}$, un mot noirci}
\[
  1 \longrightarrow V_{K} \longrightarrow T_{K} \longrightarrow \mathbb{G}_{m}
  \longrightarrow 1
\]
\emph{NB} Si $E_{K}$ unités, alors $E_{K} \cap V_{K}(\mathbb{Q})$ est d'indice
$1$ ou $2$ dans $E_{K}$, et contient $E_{1,K}$ [unités \uncertain{complètes}
y compris à l'$\infty$]

\emph{Proposition}. Soit $V'_{K}$ [$= V_{K}/R$] un groupe quotient de $V_{K}$.
Les conditions suivantes sont équivalentes :

(i) L'image de $E_{1}$ dans $V'_{K}(\mathbb{Q})$ est finie, i.e.\ $R \cap E$ est
d'indice fini dans $E$.

(ii) Il existe un « module » $\mathfrak{m}$ de $K$ tel que $T_{K} \to
T_{\mathfrak{m}} = T/\overline{E}_{\mathfrak{m}}$ est nul sur $R$, i.e.\ $T' =
T/R$ est majoré par le quotient $T_{\mathfrak{m}}$ de $T$.

(iii) $(V'_{K})_{\mathbb{R}}$ est un tore compact.

(iv) Si $L$ \struck{\ill{}} est le plus grand sous-corps Weilien de $K$, et
$L_{0} = L \cap \mathbb{R} \subset L$ (d'indice $1$ ou $2$ dans $L$)
\struck{le plus grand sous-corps tot.\ réel de $K$}, \struck{Alors l'hom.\
\ill{} (qui applique $V_{K}$ dans $V_{L}$)} $R$ \struck{est dans} [contient]
le noyau de l'\struck{épi}morphisme \struck{$T$}
\[
  V_{K} \xrightarrow{\;N_{K/L}\;} V_{L}
  \xrightarrow{\;\mathrm{can}\;} V_{L}/\operatorname{Im} V_{L_{0}} ,
\]
où $V_{L_{0}} \to V_{L}$ est induit par l'inclusion $T_{L_{0}} \hookrightarrow
T_{L}$ ; i.e.\ $V'_{K}$ « est » un quotient de $V_{L}/\operatorname{Im}
V_{L_{0}}$.

\emph{Dém.} (i) signifie qu'il existe $F \subset E_{1}$ [ss-gpe] d'indice fini
dans $E_{1}$ telle que $V'$ soit majoré par $V/\overline{F}$ ($\overline{F}$,
adhérence zariskienne de $F$) i.e.\ $T'$ majoré par $V/\overline{F}$ ; comme
les $E_{\mathfrak{m}}$ sont d'indice fini dans $E_{1}$, on voit que (ii)
$\Rightarrow$ (i), et (i) $\Rightarrow$ (ii) car tout sous-groupe
\struck{\ill{}} d'indice fini de $E_{1}$ (ou de $E$, c'est kif) contient un
groupe $E_{\mathfrak{m}}$ (Chevalley).

(i) $\Longrightarrow$ (iii) Car soit $F = R \cap E$, de sorte que $V'$ est
majoré par $V/\overline{F}$, donc $V'_{\mathbb{R}}$ majoré par
$V_{\mathbb{R}}/\overline{F}$ [où maintenant $\overline{F}$ désigne
l'adhérence zariskienne dans $V_{\mathbb{R}}$), et $V'(\mathbb{R}) =
V(\mathbb{R})/\overline{F}(\mathbb{R})$]. Or on sait que
\[
  V_{\mathbb{R}} \subset (\mathbb{R}^{*r_{1}} \times \mathbb{C}^{*r_{2}})
\]
\ill{} [\uncertain{formé des} systèmes $(\lambda_{i})(\mu_{j})$ tels que
$\prod \lambda_{i} \prod |\mu_{j}|^{2} = 1$], et $E_{1}$ (donc $F$) est un
\note{en interligne, au-dessus de la parenthèse :
$\simeq \mathbb{R}^{*(r_{1}+r_{2})} \times
(\mathbb{Z}/2\mathbb{Z})^{r_{1}} \times \mathbb{U}^{r_{2}}$, suivi d'un
passage biffé, « i.e.\ $(\lambda_{i})(\varepsilon_{j}), (u_{j})$ tels que
$\prod \lambda_{i} = 1$ » ; l'exposant de la parenthèse, et un signe après l'astérisque de
$\mathbb{R}^{*}$, sont \ill{}. La phrase
se poursuit au-delà du feuillet}

\end{document}
